%!TEX encoding = UTF-8 Unicode
%!TEX root = ../thesis.tex
\typeout{ACCESS - TEST LISTING \thepage}

\chapter{Test Code Listings} \label{chap:codelistings}

\section{Example: Verbatim Listing} \label{sec:verbatimlistings}

For now, the \LaTeX\/ package \texttt{listing} is not compatible with tagging which is why we will use the plain verbatim environment.
The first code snippet is just using a plain verbatim environment:

\begin{verbatim}
#include <stdio.h>

int main() {
    printf("Hello, World!\n");
    return 0;
}
\end{verbatim}

The second code snippet is using a custom verbatim environment with the Adobe Source Code Pro font and line numbers:

\begin{coding}
#include <stdio.h>

int main() {
    printf("Hello, World!\n");
    return 0;
}
\end{coding}



\section{Example: Verbatim combined with tcolorbox Listing} \label{sec:tcolorboxlistings}

Additionally requires the package \texttt{tcolorbox}.

\begin{tcolorbox}[colback=gray!25, boxsep=0pt, arc=0pt, boxrule=0pt]
\begin{coding}
#include <stdio.h>

int main() {
    printf("Hello, World!\n");
    return 0;
}
\end{coding}
\end{tcolorbox}

\section{Example: Code Listing with Listings}

(underlying package not compatible yet)

% \begin{lstlisting}[language=C, caption={Firmware for instruction skipping.}, label=lst:stmdut]
%     uint32_t glitched = 0x01u;
%     sigTermInit();
%     sigEnabOne();

%     while (glitched) //THIS INSTRUCTION CAN BE GLITCHED
%     {
%     ;
%     }
%     // Increment instruction skip counter
%     glitched++;
%     glitched++;
%     asm("NOP");
%     asm("NOP");
%     asm("NOP");
%     asm("NOP");
%     asm("NOP");
%     asm("NOP");

%     while (glitched)
%     {
%         if (glitched == 0x01u) {
%             sigTermOne();
%             sigEnabTwo();
%         } else if (glitched == 0x02u) {
%             sigTermOne();
%             sigTermTwo();
%         } else if (glitched == 0x03u) {
%             sigEnabOne();
%             sigEnabTwo();
%         }
%     }
% \end{lstlisting}